\documentclass[10pt,a4paper]{amsart}
\usepackage[margin=19mm]{geometry}
\usepackage{amsmath,amssymb,mathtools}
\usepackage[scr=rsfs,cal=euler]{mathalfa}
\usepackage{xfrac}
\usepackage[unicode,hidelinks,bookmarks=false]{hyperref}
\newcommand{\ie}{i.e.\ }
\newcommand{\cf}{cf.\ }
\newcommand{\resp}{respectively\ }
\newcommand{\Subsection}{\subsection}
\providecommand{\nobreakdash}{\nobreak\hbox{-}}
\setlength{\emergencystretch}{2em}
\multlinegap0pt
\usepackage{bm}
\usepackage{bbm}
\usepackage{dsfont}
\usepackage{xspace}
\usepackage{cancel}
\usepackage{mathtools}
\usepackage{amsthm}
\makeatletter
\@ifundefined{theorem}{\newtheorem{theorem}{Theorem}}{}
\makeatother
\title[A Generalized Information Bottleneck Theory of Deep Learning]{A Generalized Information Bottleneck Theory of Deep Learning}
\author{Charles Westphal, Stephen Hailes, Mirco Musolesi}
\date{Source: arXiv:2509.26327v3 (CC BY 4.0)}
\begin{document}
\maketitle

\noindent\emph{Source and licence.} A Generalized Information Bottleneck Theory of Deep Learning. Version arXiv:2509.26327v3 (CC BY 4.0), distributed under the Creative Commons Attribution 4.0 International licence, as recorded in the arXiv licence metadata for this exact version and shown on the abstract page https://arxiv.org/abs/2509.26327v3. Full rights notice: licenses/arxiv-2509.26327-CC-BY-4.0.txt.\par\smallskip

\noindent\textit{Editorial scope.} Counted unit: the Theorem whose source label is \texttt{None} in the article (source0.tex lines 365--382, source label \texttt{None}), delivered together with the article's own complete original proof of it (source0.tex lines 384--410, one \texttt{proof} environment of 27 lines). No mathematics is written, repaired or re-derived: statement and proof are the article's own text line for line. Required context retained uncounted: none. Every definition, display and label of those lines is retained unchanged. Omitted from this package and not redistributed: 1--364, 383--383, 411--742. Nothing inside a retained excerpt is omitted or altered: every retained body is delivered exactly as the article's lines are written, so no omission receipt of any kind accompanies this package. Citations: the retained excerpts carry no citation command at all, so no bibliography entry of the article is transcribed and no reference is redistributed. Reference adaptation: none was needed and none was made; every reference inside the delivered unit resolves inside it, so no unresolved reference or citation is produced. Printed slips kept literal: no printed word, symbol, hypothesis or number of the counted statement or of its proof was altered. Layout and class: the article's own class is \texttt{article} with packages including iclr2026\_conference, amsmath, amssymb, graphicx, hyperref, url, wrapfig, inputenc, fontenc, tikz, booktabs, nicefrac, microtype, tabularx; the wrapper is the lane's pinned \texttt{amsart} editorial wrapper at 10pt on A4, which restates the theorem-like environments the retained text needs and adds the article's own macro definitions that the retained text needs (none), with extra wrapper packages (bm, bbm, dsfont, xspace, cancel, mathtools). The delivered numbering is the unit's own. Assets: no retained span contains a figure environment, a graphics-inclusion command, a PDF-page inclusion, a table or a TikZ picture; no image file is copied into this package, the package declares no asset dependency and no image file is redistributed with it. Licence: version arXiv:2509.26327v3 of this article is distributed under the Creative Commons Attribution 4.0 International licence, as recorded in the arXiv licence metadata for this exact version and shown on the abstract page https://arxiv.org/abs/2509.26327v3. A full rights notice is delivered at licenses/arxiv-2509.26327-CC-BY-4.0.txt. Characters: every retained excerpt is pure ASCII, so the delivered engine, XeLaTeX with its default Latin Modern text font, reports no missing character.
\begin{theorem}
For $\beta\in(0,\infty]$, consider the GIB functional
\[
\mathcal{J}_\beta(P_{Z|X})
:= I(Z;Y)
- \frac{1}{\beta\,2N}\sum_{i=1}^N 
\Big( I(X^{-i};Q) + I(X^{i};Q) \Big).
\]
Then at $\beta=\infty$,
\[
\sup_{P_{Z|X}} \mathcal{J}_\infty
= \sup_{P_{Z|X}} I(Z;Y)
\;\le\; H(Y).
\]
Moreover, equality holds (i.e., $\sup_{P_{Z|X}} I(Z;Y)=H(Y)$) if and only if
$Y$ is a deterministic function of $Z$, equivalently $H(Y\mid Z)=0$.
In this case, $Z$ is a sufficient statistic.
\end{theorem}

\begin{proof}
At $\beta=\infty$, the penalty term vanishes and therefore
$\mathcal{J}_\infty = I(Z;Y)$.  
By the elementary information inequality
$I(Z;Y) \le H(Y)$ for all encoders $P_{Z|X}$, we have
\[
\sup_{P_{Z|X}} \mathcal{J}_\infty
= \sup_{P_{Z|X}} I(Z;Y)
\le H(Y).
\]

If for some encoder we have $H(Y\mid Z)=0$, i.e.\ $Y$ is a deterministic 
function of $Z$, then
\[
I(Z;Y) = H(Y) - H(Y\mid Z) = H(Y),
\]
and the upper bound is achieved.

Conversely, if $I(Z;Y)=H(Y)$ for some encoder, then
\[
H(Y\mid Z) = H(Y) - I(Z;Y) = 0,
\]
so $Y$ must be a deterministic function of $Z$.

Thus, at $\beta=\infty$, maximising $\mathcal{J}_\infty$ recovers precisely
the set of encoders for which $Z$ is a sufficient statistic for $Y$.
\end{proof}

\end{document}
